\section{Apéndice} 
\subsection{Expresiones teóricas de $H(\omega)$}
Los circuitos RC y RLC son sistemas LIT, por lo que si se los alimenta con una señal sinusoidal de frecuencia $\omega$ y amplitud $\text{A}$ de la forma 

\begin{equation}
    V_e=Ae^{i\omega t}
    \label{eq:1}
\end{equation}
la señal de salida será \cite{sistemasLIT}
\begin{equation}
    V_S=H(\omega)V_E
    \label{eq:2}
\end{equation}
donde la función de transferencia dependerá de la componente sobre la que se mida. 
    Aplicando la Ley de Kirchoff  para el circuito RC, se puede obtener la ecuación diferencial para el voltaje sobre la resistencia o sobre el capacitor\cite{kip1969}. Se obtiene entonces para el voltaje de salida de la resistencia
\begin{equation}
    V_e=V_R+\frac{1}{\tau}\int V_Rdt
\end{equation}
siendo $\tau=RC$ la constante de tiempo del circuito. Aplicando lo conocido sobre sistemas LIT se obtiene
\begin{equation}
    V_e=H(\omega)V_e + \frac{H(\omega)}{\tau}\int V_edt
\end{equation}
donde $H(\omega)$ es constante al ser el sistema invariante en el tiempo. Integrando $V_e$ se llega a la expresión
\begin{equation}
    V_e=H(\omega)V_e+\frac{H(\omega)V_e}{i\omega\tau}
\end{equation}
o lo que es equivalente
\begin{equation}
    H(\omega)(1+\frac{1}{i\omega\tau})=1
\end{equation}
y por lo tanto se tiene que el módulo de la función de transferencia es
\begin{equation}
    |H_r(\omega)|=\frac{\omega\tau}{\sqrt{1+(\omega\tau)^2}}
    \label{eq:moduloRC}
\end{equation}
y su desfasaje con respecto al voltaje de entrada
    \begin{equation}
    \varphi_r=\arctan(\frac{1}{\omega\tau})
    \label{eq:3}
\end{equation}
De la misma manera se puede demostrar para el capacitor el modulo de la función de transferencia es
\begin{equation}
    |H_C(\omega)|=\frac{1}{\sqrt{1+(\omega\tau)^2}}
\end{equation}
y el desfasaje entre el voltaje de entrada y el del capacitor es
\begin{equation}
    \varphi_c=\arctan(- \omega\tau)
    \label{eq:4}
\end{equation}
    Por otro lado, si se tiene circuito RLC en serie (figura \ref{fig:RLC}) y se le 
aplica un voltaje de entrada $V_e$, la ecuación de la ley de Kirchoff queda
\begin{equation}
    V_e=V_R+\frac{1}{\tau\omega_0^2}\frac{dV_R}{dt}+\frac{1}{\tau}\int V_Rdt
\end{equation}
    donde se definió $\omega_0=\frac{1}{\sqrt{LC}}$ la frecuencia resonante del sistema. Aplicando los mismos conceptos que para el circuito RC, se obtiene que el módulo de la función de transferencia, cuando se mide sobre la resistencia, es
\begin{equation}
    |H(\omega)|= \frac{\omega\tau}{\sqrt{(\omega\tau)^2+(1-(\frac{\omega}{\omega_0})^2)^2}}
    \label{eq:moduloRLC}
\end{equation}
y el desfasaje con respecto al voltaje de entrada
\begin{equation}
    \tan(\varphi)=\frac{(\frac{\omega}{\omega_0})^2-1}{\omega\tau}
    \label{eq:desfasajeRLC}
\end{equation}
    A su vez, al ser el voltaje de la resistencia directamente proporcional a la corriente del circuito, es posible hallar el módulo y el argumento de la impedancia del sistema a partir de la relación
    \begin{equation}
        V_R=IR \ \ \ \ \ \ V_e=IZ(\omega) \Longrightarrow V_R=\frac{V_eR}{|Z(\omega)|}e^{-i\varphi}
    \end{equation}
    donde el argumento $-\varphi$ debe ser igual al de $H(\omega)$ anteriormente encontrado, ya que el voltaje de la resistencia se encuentra en fase con la corriente del sistema. Finalmente $\frac{R}{|Z(\omega)|}$ debe ser igual a $|H(\omega)|$ por igualdad de módulos de números complejos, obteniendo así que el módulo de $Z(\omega)$ es
    \begin{equation}
        |Z(\omega)|=\frac{R}{\omega\tau} \sqrt{(\omega\tau)^2+\left(1-\left(\frac{\omega}{\omega_0}\right)^2\right)^2}
        \label{imp RCL}
    \end{equation}
      



